
\section{The seven-line normal form}\label{sec:normalform}

Houston~\cite{HOU16A,HOU16B}, starting from the classical Lucas
parametrization, explained this almost-magic normal form and its
interpretation as three three-term arithmetic progressions of one common
difference. We retain a self-contained proof because the integral image
lattice below is needed to identify the exact congruence obstruction.

\begin{lemma}[Retraction bookkeeping]\label{lem:retraction}
Let a fully magic $3\times3$ array have common line sum $M$.
\emph{(i)} Subtracting the same $t$ from the three cells of one long
diagonal leaves the three rows, the three columns and the other diagonal
with sum $M-t$, while the selected diagonal has sum $M-3t$; the exact line
defect is $2t$.
\emph{(ii)} More generally, subtracting $\tau$ from the three cells of a
permutation transversal $T$ leaves every row and every column with sum
$M-\tau$; if $T$ meets the two diagonals in $j$ and $k$ cells, their sums
become $M-j\tau$ and $M-k\tau$.
\emph{(iii)} Conversely, if exactly one diagonal of a $3\times3$ array
fails while the other seven lines share a common sum, adding half the
defect to each cell of the failed diagonal restores a fully magic array.
\end{lemma}

\begin{proof}
(i) Every row and every column contains exactly one modified cell, and the
other diagonal contains exactly one modified cell, while the selected
diagonal contains all three; the sums follow, and the defect is
$(M-t)-(M-3t)=2t$. Nothing here uses squareness or any property of a
particular family.
(ii) A permutation transversal meets every row and every column in exactly
one cell, and meets the diagonals in $j$ resp.\ $k$ cells by definition of
$j,k$; the same count gives the sums.
(iii) By Theorem~\ref{main:normalform} (proved independently below) the
array is the displayed normal form, with seven-line sum $A+E+P$, failed
antidiagonal sum $3E$, and defect $\delta=A+P-2E$. Each row and column
contains exactly one antidiagonal cell, and the principal diagonal contains
one (the centre), so adding $\delta/2$ to the three antidiagonal cells
raises all seven line sums to $A+E+P+\delta/2$ and the antidiagonal to
$3E+3\delta/2$; both equal $\tfrac32(A+P)$, so all eight lines agree.
\end{proof}

\begin{proposition}[Integral line-sum lattice]\label{prop:linesumsnf}
Let $\Lambda:\mathbb Z^9\to\mathbb Z^8$ send
$X=(a,b,c;d,e,f;g,h,i)$ to
$(R_1,R_2,R_3,C_1,C_2,C_3,D_+,D_-)$, its three row sums, three column sums
and two diagonal sums. Then
\[
\operatorname{SNF}(\Lambda)=\operatorname{diag}(1,1,1,1,1,1,3,0),
\qquad
\operatorname{coker}(\Lambda)\cong\mathbb Z\oplus\mathbb Z/3\mathbb Z.
\]
An integer vector of eight prescribed line sums lies in the image if and
only if
\[
R_1+R_2+R_3=C_1+C_2+C_3,
\qquad
R_1+R_3\equiv C_2+D_++D_-\pmod3.
\]
Consequently $\Lambda$ has rank $6$ over $\mathbb F_3$ and rank $7$ over
$\mathbb F_p$ for every prime $p\ne3$.
\end{proposition}

\begin{proof}
In the entry order $(a,b,c,d,e,f,g,h,i)$, the matrix of $\Lambda$ is
\[
\left(\begin{array}{ccccccccc}
1&1&1&0&0&0&0&0&0\\0&0&0&1&1&1&0&0&0\\
0&0&0&0&0&0&1&1&1\\1&0&0&1&0&0&1&0&0\\
0&1&0&0&1&0&0&1&0\\0&0&1&0&0&1&0&0&1\\
1&0&0&0&1&0&0&0&1\\0&0&1&0&1&0&1&0&0
\end{array}\right).
\]
The primitive left relation
$-R_1-R_2-R_3+C_1+C_2+C_3=0$ bounds the rank by $7$. The minor on rows
$(R_1,R_2,R_3,C_1,C_2,D_+,D_-)$ and columns $(a,b,c,d,e,f,g)$ is $-3$,
whereas the minor on rows $(R_1,R_2,R_3,C_1,C_2,D_+)$ and columns
$(a,b,c,d,f,g)$ is $-1$. Moreover
\[
R_1+R_3-C_2-D_+-D_-=-3e,
\]
so modulo $3$ the rank is at most $6$ and every $7$-minor is divisible by
$3$. The sixth and seventh determinantal divisors are therefore $1$ and
$3$, giving the displayed Smith form. The two stated conditions annihilate
the image and define a surjective map
$\mathbb Z^8\to\mathbb Z\oplus\mathbb Z/3\mathbb Z$ (use the $R_2$ and
$D_+$ coordinate vectors); its kernel has the same quotient as
$\operatorname{im}\Lambda$, so the two coincide. Reduction of the Smith
form gives the rank statement.
\end{proof}

\begin{proof}[Proof of Theorem~\ref{main:normalform}]
The failed line cannot be a row: the three columns already force the total
$3S$, so if two rows have sum $S$ the third does as well. The same argument
excludes a failed column. Up to the dihedral action of the grid the failed
line is therefore the antidiagonal. Proposition~\ref{prop:linesumsnf}
forces its sum to be divisible by $3$. Solving the seven linear line
equations then shows that its central cell is one third of that sum and
gives the displayed universal form; the seven common lines have sum
$A+E+P$ and the antidiagonal has sum $3E$. The nine entries decompose into
the three arithmetic progressions
\[
(A-D,\,A,\,A+D),\qquad (P-D,\,P,\,P+D),\qquad (E-D,\,E,\,E+D),
\]
of one common difference $D$. Consequently an array with exactly eight
square entries consists of two full three-square progressions and one
progression containing exactly two squares, which is the last assertion.
\end{proof}

\begin{remark}[Where full magic lives]\label{rem:fourthAP}
In the normal form, the array is fully magic if and only if the
antidiagonal sum equals the common sum, that is $3E=A+E+P$, equivalently
\[
A+P=2E:
\]
the three centres themselves must form a three-term arithmetic progression.
For an all-squares array this is a \emph{fourth} square progression, with
its own common difference $E-A$ --- never equal to $\pm D$, since $E=A+D$
would repeat the centre at positions $(2,2)$ and $(2,3)$, while $E=A-D$
would repeat it at $(2,2)$ and $(3,2)$, against nine-way distinctness ---
and this single extra condition is exactly
where the elliptic wall of Section~\ref{sec:surface} lives. The classical
one-parameter families of nine-square arrays with seven equal sums
(Lucas~\cite{LUC76}; Sallows's record example~\cite{SAL97}) sit inside this
normal form, as does the general parametrization of $3\times3$ magic
squares going back to Lucas~\cite{LUC91}.
\end{remark}

\begin{remark}[Census normalization]\label{rem:census}
The bounded census of Appendix~A quotients by the geometric $D_4$ action
and measures height by the largest square root among the eight square
entries; its exact minimum statements are supporting propositions at their
frozen bound and are never used as evidence for an unbounded claim.
\end{remark}
